\documentclass[11pt,twoside]{amsart}
\usepackage{amssymb, amsmath, enumerate, palatino, hyperref}
\usepackage{caption}
\usepackage[normalem]{ulem}
\usepackage{fullpage}
\usepackage{tikz}
\usetikzlibrary{arrows.meta,positioning}

\usepackage[T1]{fontenc}
\renewcommand{\labelitemi}{\guillemotright}
\usepackage{mathrsfs}
\hypersetup{hidelinks}

\theoremstyle{plain}
\newtheorem{prop}{Proposition}
\newtheorem{lemma}[prop]{Lemma}
\newtheorem{thm}[prop]{Theorem}
\newtheorem{cor}[prop]{Corollary}
\theoremstyle{remark}
\newtheorem{rmk}[prop]{Remark}
\newtheorem*{rmk*}{Remark}
\newtheorem{prob}{Problem}
\newtheorem*{bonus*}{Bonus Problem}
\theoremstyle{definition}
\newtheorem{ex}[prop]{Example}
\newtheorem{defn}[prop]{Definition}

\newcommand{\RR}{\mathbb{R}}
\newcommand{\ZZ}{\mathbb{Z}}
\newcommand{\Ztn}{(\mathbb{Z}/2)^{n}}
\newcommand{\tr}{\operatorname{tr}}
\newcommand{\spec}{\operatorname{spec}}
\newcommand{\wt}{\omega}

\title{Math 372: Combinatorics\\ Homework 02}
\author{Due Friday, 25 September, 10\textsc{p.m.}, via Gradescope}

\begin{document}
\maketitle

\begin{rmk*}
\textbf{When you can start.} Problems 1 and 2 use only what we have done through Monday,
September~14. Problem~3 needs Wednesday the 16th, and Problem~4 needs Friday the 18th. The set is
released early so that the first half can be done while the lab is also open; do not wait for
Friday to begin.
\end{rmk*}

\begin{rmk*}
\textbf{Conventions.} Graphs are finite and undirected. For $u,v\in\Ztn$ the dot product
$u\cdot v=\sum_i u_iv_i$ is taken mod~$2$, $\wt(u)$ is the number of $1$'s in $u$, and
$\chi_u\colon\Ztn\to\RR$ is the function $\chi_u(v)=(-1)^{u\cdot v}$, as in class. A poset is
finite; $P_i$ denotes the $i$th level of a graded poset and $p_i=\#P_i$. An \emph{order-matching}
from $P_i$ to $P_{i+1}$ is a one-to-one map $\mu\colon P_i\to P_{i+1}$ with $x<\mu(x)$ for all
$x$; from $P_{i}$ to $P_{i-1}$ it is a one-to-one map with $\mu(x)<x$.
\end{rmk*}

\begin{rmk*}
\textbf{On these solutions.} Each problem is worth up to five points for mathematical content
and up to two for writing. Write for another student in this course: say what you are going to
show before you show it, justify each step, and use complete sentences. End your submission with
a short \emph{Methods and tools} note --- three or four sentences on what software you used and
for what, and how you checked anything a machine produced for you.
\end{rmk*}

%%% Problem 1 %%%
\begin{prob}[the two-flip graph]
Fix $n\ge3$. Let $G_n$ be the graph with vertex set $\Ztn$ in which $u$ and $v$ are adjacent
exactly when they differ in \emph{exactly two} coordinates. (So $G_n$ is to the cube as ``change
two bits'' is to ``change one bit.'')
\begin{enumerate}[(a)]
\item Show that $G_n$ is regular and find its degree.
\item Show that each $\chi_u$ is an eigenvector of the adjacency matrix of $G_n$, and find its
eigenvalue as a function of $\wt(u)$. Then write out $\spec(G_4)$ completely, with
multiplicities.
\item How many connected components does $G_n$ have? Describe them. (Which vertices can a walk
starting at $0$ reach?)
\item Use $\spec(G_4)$ to compute the number of closed walks of length $3$ in $G_4$, and deduce
how many triangles $G_4$ contains. Then confirm that number by a direct count, using your
description of the components in (c).
\end{enumerate}
\end{prob}

%%% Problem 2 %%%
\begin{prob}[coming home]
Let $G$ be a $d$-regular graph on $p$ vertices with adjacency matrix $A$, and let $M=\tfrac1dA$,
so that $(M^{\ell})_{uv}$ is the probability that the simple random walk started at $u$ is at
$v$ after $\ell$ steps. Write $p_\ell(u)=(M^{\ell})_{uu}$ for the probability of being home after
$\ell$ steps.
\begin{enumerate}[(a)]
\item Explain why $p_\ell(u)$ does not depend on $u$ when $G$ is the complete graph $K_p$, and
also when $G$ is the cube $Q_n$. Conclude that in both cases $p_\ell=\tr(M^{\ell})/p$.
\item For $K_p$, find $p_\ell$ exactly as a function of $p$ and $\ell$, and compute
$\lim_{\ell\to\infty}p_\ell$.
\item For $Q_n$, show that $p_\ell=0$ when $\ell$ is odd, find $\lim_{m\to\infty}p_{2m}$, and
explain in one sentence why the answer is $2^{1-n}$ rather than $2^{-n}$.
\end{enumerate}
\end{prob}

%%% Problem 3 %%%
\begin{prob}[a poset where the matchings can be written down]
Fix integers $a,b\ge1$ and let $P$ be the set of pairs $(i,j)$ with $0\le i\le a$ and
$0\le j\le b$, ordered by $(i,j)\le(i',j')$ if and only if $i\le i'$ and $j\le j'$. (If you like:
$P$ is the set of divisors of $2^a3^b$, ordered by divisibility.)
\begin{enumerate}[(a)]
\item Show that $P$ is graded of rank $a+b$, with $P_k=\{(i,j):i+j=k\}$. Compute the rank sizes
$p_0,\dots,p_{a+b}$ and show the sequence is symmetric and unimodal.
\item Assume $a\le b$. Find order-matchings
\[
P_0\to P_1\to\cdots\to P_b \qquad\text{and}\qquad P_b\leftarrow P_{b+1}\leftarrow\cdots\leftarrow P_{a+b},
\]
and verify that they are order-matchings. (There is a very simple choice; try changing only one
coordinate.)
\item Conclude that the largest antichain in $P$ has exactly $\min(a,b)+1$ elements.
\item Draw the chain decomposition of $P$ that your matchings produce when $a=2$ and $b=3$.
\end{enumerate}
\end{prob}

%%% Problem 4 %%%
\begin{prob}[the linear algebra method]
Recall the operators on the Boolean algebra $B_n$: for a set $x$ with $|x|=i$,
\[
U_i(x)=\sum_{\substack{|y|=i+1\\ y\supset x}}y,
\qquad
D_i(x)=\sum_{\substack{|z|=i-1\\ z\subset x}}z,
\]
extended linearly to $U_i\colon\RR(B_n)_i\to\RR(B_n)_{i+1}$ and
$D_i\colon\RR(B_n)_i\to\RR(B_n)_{i-1}$.
\begin{enumerate}[(a)]
\item For $n=3$, write out the matrices $[U_1]$ and $[D_2]$ with respect to the bases
$(B_3)_1=\{1,2,3\}$ and $(B_3)_2=\{12,13,23\}$, in those orders. Then compute
$[D_2][U_1]-[U_0][D_1]$ and check that it is what the identity from class predicts.
\item Now let $P$ be any graded poset and suppose $U\colon\RR P_i\to\RR P_{i+1}$ is linear,
\emph{order-raising} (each $U(x)$ is a linear combination of elements $y>x$), and \emph{onto}.
Prove that there is an order-matching from $P_{i+1}$ to $P_i$.

In class we proved the companion statement --- \emph{one-to-one} order-raising gives a matching
from $P_i$ to $P_{i+1}$ --- by extracting a nonzero term from a determinant. Your proof should
run the same way, but be careful about which rows and columns you are choosing and why.
\end{enumerate}
\end{prob}

%%% Bonus %%%
\begin{bonus*}[LYM]
Let $\mathcal A$ be an antichain in $B_n$. Show that
\[
\sum_{A\in\mathcal A}\binom{n}{|A|}^{-1}\ \le\ 1,
\]
and deduce Sperner's theorem from it. (Count, in two ways, the pairs $(\mathcal C,A)$ where
$\mathcal C$ is a maximal chain $\emptyset\subset\{x_1\}\subset\{x_1,x_2\}\subset\cdots$ in $B_n$
and $A\in\mathcal A$ lies on $\mathcal C$. How many maximal chains pass through a given $A$? How
many can pass through two different members of an antichain?)
\end{bonus*}

\end{document}
